Sigui el sistema d'equacions lineals
\[
\left.\begin{aligned}
2x+y&=1+z\\
my+z&=2-x\\
mz+3&=3x+y
\end{aligned}\right\},
\]
en què $m$ és un nombre real.

\begin{apartats}

\apartat{1,25}
Discutiu el sistema segons els valors del paràmetre $m$.

\begin{solucio}
El sistema, en la forma normal, és
$\left\{\begin{aligned}2x+y-z&=1\\x+my+z&=2\\3x+y-mz&=3\end{aligned}\right.$.
Calculem el determinant de la matriu dels coeficients:
\[
|A|=\begin{vmatrix}2&1&-1\\1&m&1\\3&1&-m\end{vmatrix}=\left(-2m^2-1+3\right)-(-3m-m+2)=-2m^2+4m=2m(2-m).
\]
$|A|=0\rightarrow m=0$ i $m=2$, que són els valors que cal tenir en compte en la discussió. Tenim tres casos:\\
Si $m\neq0$ i $m\neq2$, com que $\operatorname{rang}(A)=3$ [$|A|\neq0$] $=\operatorname{rang}(A^*)$ $=$ nombre
d'incògnites, tenim un \textbf{SCD} (sistema compatible determinat): \textbf{solució única}.\\
Si $m=0$, per Gauss obtenim
\[
A^*=\left(\begin{array}{ccc|c}2&1&-1&1\\1&0&1&2\\3&1&0&3\end{array}\right)
\sim\left(\begin{array}{ccc|c}1&0&1&2\\2&1&-1&1\\3&1&0&3\end{array}\right)
\sim\left(\begin{array}{ccc|c}1&0&1&2\\0&-1&3&3\\0&-1&3&3\end{array}\right)
\sim\left(\begin{array}{ccc|c}1&0&1&2\\0&-1&3&3\\0&0&0&0\end{array}\right)
\]
(intercanviant $f_1$ i $f_2$; fent $2f_1-f_2$ i $3f_1-f_3$; i fent $f_3-f_2$).
$\operatorname{rang}(A)=2=\operatorname{rang}(A^*)<3=$ nombre d'incògnites. Per tant, es tracta d'un \textbf{SCI}
(sistema compatible indeterminat) amb $3-2=1$ grau de llibertat: infinites solucions, que depenen d'un paràmetre.\\
Si $m=2$, per Gauss obtenim
\[
\left(\begin{array}{ccc|c}2&1&-1&1\\1&2&1&2\\3&1&-2&3\end{array}\right)
\sim\left(\begin{array}{ccc|c}1&2&1&2\\2&1&-1&1\\3&1&-2&3\end{array}\right)
\sim\left(\begin{array}{ccc|c}1&2&1&2\\0&3&3&3\\0&5&5&3\end{array}\right)
\sim\left(\begin{array}{ccc|c}1&2&1&2\\0&1&1&1\\0&0&0&-6\end{array}\right)
\]
(intercanviant $f_1$ i $f_2$; fent $2f_1-f_2$ i $3f_1-f_3$; i fent $f_2/3$ i $3f_3-5f_2$).
$\operatorname{rang}(A)=2\neq\operatorname{rang}(A^*)=3\Rightarrow$ \fbox{\textbf{SI} (sistema incompatible): no hi ha
solució.}

\textit{Pauta oficial:} 0,25 pel plantejament matricial i el determinant de la matriu de coeficients, 0,25 pel
càlcul dels valors singulars, i 0,75 per la discussió (0,25 per cada cas). L'apartat es pot resoldre també via el
càlcul de rangs.
\end{solucio}

\apartat{1,25}
Resoleu el sistema, si té solució, per al cas $m=1$.

\begin{solucio}
Si $m=1$, quan substituïm en el sistema obtenim
\[
A^*=\left(\begin{array}{ccc|c}2&1&-1&1\\1&1&1&2\\3&1&-1&3\end{array}\right)
\sim\left(\begin{array}{ccc|c}1&1&1&2\\2&1&-1&1\\3&1&-1&3\end{array}\right)
\sim\left(\begin{array}{ccc|c}1&1&1&2\\0&1&3&3\\0&2&4&3\end{array}\right)
\sim\left(\begin{array}{ccc|c}1&1&1&2\\0&1&3&3\\0&0&-2&-3\end{array}\right)
\]
(intercanviant $f_1$ i $f_2$; fent $2f_1-f_2$ i $3f_1-f_3$; i fent $f_3-2f_2$), és a dir,
\[
\left\{\begin{aligned}x+y+z&=2\\y+3z&=3\\2z&=3\end{aligned}\right.
\;\rightarrow\;
\left\{\begin{aligned}x-\tfrac32+\tfrac32&=2\\y+3\cdot\tfrac32&=3\\z&=\tfrac32\end{aligned}\right.
\;\rightarrow\;
\boxed{\;x=2,\quad y=-\tfrac32,\quad z=\tfrac32\;}
\]

\textit{Pauta oficial:} 0,25 per la substitució del paràmetre, 0,25 pel procediment de resolució, i 0,75 per la
resolució (0,25 per cada variable). L'apartat es pot resoldre també pel mètode de Cramer.
\end{solucio}

\end{apartats}
